\chapter{NLF en detalle}
\label{cap:nlf-detalle}

\begin{intuicion}
Capítulo~\ref{cap:prune-nlf} introdujo NLF como prune; este
capítulo muestra cómo está implementado por dentro: layout
\texttt{[u32; 9]} por vértice, \texttt{IndexBacking} para la
dualidad heap/mmap, builder \(O(|\E|)\) en una pasada, e
incremental update detallado más allá del overview de
\ref{cap:incremental-nlf}. Es el primer índice ``simple'' del
motor; entender su implementación a fondo es buena preparación para
el más complejo (k-hop, próximo capítulo).
\end{intuicion}

\section{Layout}

\begin{lstlisting}[language=Rust]
const NLF_TYPE_COUNT: usize = 9;  // tag de EntityType en [0, 9)

pub struct NlfTable {
    counts: IndexBacking<[u32; NLF_TYPE_COUNT]>,
}
\end{lstlisting}

\noindent
La estructura es trivial: una tabla con \(|\V|\) filas, cada fila
es \texttt{[u32; 9]} (36 bytes). El \texttt{IndexBacking} (visto
en el capítulo anterior) decide si vive en heap o en mmap.

\subsection{Por qué \texttt{[u32; 9]} y no \texttt{HashMap<EntityType, u32>}}

Un \texttt{HashMap} sería más flexible pero terrible en
performance. \(N \cdot \texttt{HashMap}\) implica:

\begin{itemize}
    \item Por cada lookup, hash + bucket walk + compare. \(\sim 30\)
    ciclos.
    \item Memoria: $\sim$48 B por bucket + factor de carga
    (\emph{load factor}) típico 0.7. Para 9 entries por vértice:
    $\sim$72 B/vértice. Versus 36 B del array.
    \item Locality: los buckets están dispersos en heap; iterar
    el HashMap toca cache lines aleatorias.
\end{itemize}

\noindent
El array fijo \texttt{[u32; 9]}:

\begin{itemize}
    \item Lookup: indexación directa, 1 ciclo.
    \item Memoria: 36 B per row, contiguous.
    \item Locality: una fila entera entra en una cache line (con
    20 B de margen para los siguientes 5 rows).
\end{itemize}

\noindent
Trade-off: el array fija 9 tipos. Si en el futuro se agrega un
décimo \texttt{EntityType}, hay que crecer el array (cambio
breaking de schema). El equipo decidió que la fijación de tipos es
un commitment razonable.

\section{El builder}

\begin{algorithm}[H]
\caption{\texttt{NlfTable::build}}
\KwIn{grafo \(G\) (sealed)}
\KwOut{\texttt{NlfTable}}
\(\mathrm{counts} \leftarrow \mathrm{IndexBacking::heap\_zeros}(|\V|)\)\;
\For{cada \texttt{cr} en \texttt{G.edge\_store.iter\_all()}}{
    \(\mathrm{src} \leftarrow \mathrm{cr.source\_sid.index()}\)\;
    \(L \leftarrow \mathrm{cr.dest\_type\_tag()}\)\;
    \If{\(L < \texttt{NLF\_TYPE\_COUNT}\)}{
        \(\mathrm{counts}[\mathrm{src}][L] \leftarrow \mathrm{counts}[\mathrm{src}][L].\mathrm{saturating\_add}(1)\)\;
    }
}
\Return \texttt{NlfTable\{ counts \}}\;
\end{algorithm}

\noindent
Una pasada secuencial sobre \texttt{edge\_store} (que ya está
sorted, capítulo~\ref{cap:finalize}). Cada iteración: índice
directo + saturating add. \(O(|\E|)\) total.

Costo en hardware: para \(|\E| = 10^8\), \(\sim 30\) segundos
(memoria-bound, $\sim$3 MB/s de actualizaciones). El \texttt{lazy
build} (capítulo~\ref{cap:finalize}) lo difiere al primer hunt
que consulta NLF.

\subsection{¿Por qué \texttt{dest\_type\_tag} y no \texttt{source\_type\_tag}?}

NLF cuenta \emph{vecinos salientes} de un vértice por tipo destino.
La fila se indexa por source (de la arista), pero la columna se
indexa por tipo del destino. Es lo que el prune del primer paso
necesita: ``¿\(s\) tiene vecinos del tipo que el primer paso
exige?''.

Una variante alternativa: NLF \emph{simétrica} que indexa por
ambos pares. El motor no la implementa porque las hipótesis suelen
ser direccionadas (origen $\rightarrow$ destino) y la simetría
duplicaría el costo sin beneficio claro.

\section{Lookup: \texttt{can\_match}}

\begin{lstlisting}[language=Rust]
impl NlfTable {
    pub fn can_match(
        &self,
        sid: StrId,
        required: &[(u8, u32)],   // (dest_type_tag, min_count)
    ) -> bool {
        let row = self.counts.get(sid.index());
        let Some(row) = row else { return false };  // sid out of range
        required.iter().all(|&(tag, min)| {
            tag as usize >= NLF_TYPE_COUNT || row[tag as usize] >= min
        })
    }
}
\end{lstlisting}

\noindent
\(O(|\text{required}|)\) que es \(\leq |\Lc_V| = 9\). Para
hipótesis chain (un solo destination type), es \(O(1)\).

\subsection{Cómo lo invoca el matcher}

El DFS llama \texttt{can\_match} para cada candidate start:

\begin{lstlisting}[language=Rust]
let nlf = graph.ensure_nlf();
let required = collect_required_destinations(&hypothesis);
for sid in candidate_starts {
    if !nlf.can_match(sid, &required) {
        continue;  // skip this start
    }
    // ... DFS desde sid
}
\end{lstlisting}

\noindent
Para grafos con \(|\V_{L_0^V}| = 10^5\) candidates, son \(10^5\)
calls a \texttt{can\_match}: $\sim$1 ms total. El ahorro en DFS
descartado suele ser de varios segundos.

\section{Memoria}

\begin{itemize}
    \item Heap-resident: \(|\V| \cdot 36\)~B. Para \(|\V| = 10^7\):
    360 MB. Para \(10^8\): 3.6 GB.
    \item mmap-spilled: misma memoria total, pero residente en
    disco con paginación on-demand.
    \item Budget default: 1 GB. Triggeará spill alrededor de \(|\V|
    = 28 \cdot 10^6\).
\end{itemize}

\subsection{Cuándo spillea}

\begin{lstlisting}[language=Rust]
const NLF_BUDGET_BYTES: usize = 1 << 30;  // 1 GB

pub fn build(graph: &GraphHunter) -> Result<Self, Error> {
    let n_vertices = graph.vertex_count();
    let counts = IndexBacking::build_with_budget(
        n_vertices,
        NLF_BUDGET_BYTES,
        |v| populate_counts(v, graph),
    )?;
    Ok(Self { counts })
}
\end{lstlisting}

\noindent
\texttt{build\_with\_budget} (capítulo~\ref{cap:indices-auxiliares})
decide al final del populate: si los bytes acumulados exceden el
budget, vuelca a archivo y mmapea read-only.

\section{Incremental update revisitado}

Capítulo~\ref{cap:incremental-nlf} introdujo \texttt{try\_record\_edge}:

\begin{lstlisting}[language=Rust]
impl NlfTable {
    pub fn try_record_edge(
        &self,
        src_sid: StrId,
        dst_tag: u8,
    ) -> bool {
        if dst_tag as usize >= NLF_TYPE_COUNT {
            return true;  // no-op
        }
        let slice = self.counts.heap_slice_mut();
        let Some(slice) = slice else {
            return false;  // mmap-spilled
        };
        if let Some(row) = slice.get_mut(src_sid.index()) {
            row[dst_tag as usize] = row[dst_tag as usize].saturating_add(1);
        }
        true
    }
}
\end{lstlisting}

\noindent
La \emph{interior mutability} ---el método toma \texttt{\&self} no
\texttt{\&mut self}--- requiere que \texttt{IndexBacking::heap\_slice\_mut}
también la soporte. La implementación real usa
\texttt{UnsafeCell<Vec<...>>} con un invariante:
``\texttt{try\_record\_edge} se llama sólo desde el ingestor, que
es single-threaded; los matchers concurrentes leen vía
\texttt{counts.get()} que es read-only''.

\begin{advertencia}
Este patrón \emph{interior mutability single-writer} es seguro pero
delicado. Si en el futuro se permite ingest multi-threaded, el
invariante se rompe y hay que envolver en un \texttt{RwLock}. El
test de \texttt{Loom} (model checker para concurrencia)
verifica el invariante actual.
\end{advertencia}

\section{¿Y si en el futuro hace falta más resolución?}

NLF actual cuenta cardinalidad por tipo destino. No considera:

\begin{itemize}
    \item \textbf{Tipo de arista}. Si el primer paso es
    \texttt{Auth}, NLF no distingue entre vecinos vía Auth o vía
    Spawn. Si quisiéramos distinguir, sería ``NLP''
    (\emph{Neighborhood Label Profile}): \texttt{[u32; 9 \(\times\) |RelType|]}
    por vértice. Memoria \(\sim 15\times\) más.
    \item \textbf{Ventana temporal}. NLF cuenta todos los vecinos
    independiente del tiempo. Para descartar starts cuyas aristas
    están todas fuera de la ventana, el k-hop temporal sirve
    (capítulo~\ref{cap:khop-detalle}).
    \item \textbf{Doble-hop counting}. ``Cuántos vecinos-de-vecinos
    de tipo X tiene este vértice''. Sería caro de mantener
    (\(O(|\V| \cdot \deg)\) memoria); k-hop bit aproxima la
    pregunta más débil ``¿hay alguno?''.
\end{itemize}

\noindent
Cada extensión es un trade-off memoria-versus-discriminación. La
versión actual gana en simplicidad y resuelve el caso del 80\% de
los hunts.

\section{Resumen}

\begin{itemize}
    \item NLF: \texttt{[u32; 9]} por vértice. 36 B/row. Index
    directo, 1 ciclo de lookup.
    \item Builder: una pasada \(O(|\E|)\) sobre edge\_store.
    Saturating add para evitar overflow silencioso.
    \item Lookup \texttt{can\_match}: \(O(|\text{required}|)\),
    típico $\leq$ 9 ops.
    \item Memoria: 36 B \(\cdot |\V|\); spillea a mmap arriba de
    $\sim$28M vertices con budget 1 GB.
    \item Incremental update \texttt{try\_record\_edge}: interior
    mutability single-writer, atomic-free porque el invariante de
    GraphHunter lo permite.
\end{itemize}

\begin{ejercicios}
\begin{enumerate}
    \item \starred{} Calculá la memoria de NLF para un grafo
    \(|\V| = 5 \cdot 10^7\). ¿Spillea con budget 1 GB?

    \item ¿Por qué \texttt{can\_match} usa \texttt{.all()} en vez
    de \texttt{.any()}? Pensá: ¿qué quiere chequear el matcher en
    el caso multi-required?

    \item \starred{} Diseñá una variante NLP (Neighborhood Label
    Profile) que distinga por tipo de arista. ¿Cuánta memoria? ¿Qué
    cambios al builder?

    \item Si el motor cambia el patrón de ingest a multi-thread,
    ¿qué se rompe en \texttt{try\_record\_edge}? Diseñá la versión
    safe usando \texttt{AtomicU32} ---¿qué se gana, qué se pierde?

    \item Imaginá que el budget de NLF baja a 100 MB (laptop con
    poca RAM). ¿A partir de qué $|\V|$ spillea? ¿Cómo afecta el
    rendimiento del prune?
\end{enumerate}
\end{ejercicios}
